\enlargethispage{4pt}
\section{Verification}\label{sec:verification}

Every numerical assertion in this paper is checked by exact arithmetic.%
\footnote{Exact arithmetic is not a scruple here. Several of the quantities are differences of large binomial coefficients, and at $n=7$ the semiregularity target exceeds $10^{11}$; in double precision the comparison in item (V) would be meaningless. Everything uses Python's \texttt{fractions} module or integers, and nothing converts to a float.} The
scripts are described in \Cref{app:scripts}. They use the standard library of
Python, together with \texttt{sympy} for exact symbolic linear algebra in
several of them and \texttt{numpy} for integer arrays in two; they work over $\QQ$ and
its finite extensions with the \texttt{fractions} module or exact symbolic
arithmetic, and they perform no floating-point operation; items (XXIX), (XXXI) and (XL) are $\Ext$
computations over a polynomial ring, carried out in Macaulay2. We describe what is checked.

\begin{enumerate}[label=\textup{(\Roman*)},leftmargin=*]
\item \emph{The Fourier--Mukai transform on powers of the polarisation.}
  The exterior algebra $\bigwedge^{*}(V\oplus V^{\vee})$ of
  \Cref{not:model} is constructed explicitly for $g\le4$, with basis elements
  indexed by bit masks and structure constants computed from the shuffle sign.
  The class $\ell$ of \eqref{eq:ell}, the forms $\theta$ and $\theta'$ of
  \eqref{eq:thetamodel} and \eqref{eq:thetapmodel}, and the exponential
  $e^{\ell}$ are formed, and $\Phi_{\cP}(\theta'^{\,k})$ is computed by
  extracting the coefficient of the top form in the dual variables. The output
  is checked to be an exact rational multiple of $\theta^{\,g-k}$ with no
  other component, and the multiple is compared with \eqref{eq:fmtheta}. All
  values agree.

\item \emph{The sign.} The uniformity of the sign $\sigma(g,k)$ over the
  index sets $T$, which is the step in the proof of \Cref{thm:fmtheta} that
  the reader may wish to see confirmed, is checked as a permutation
  computation for every $g\le12$ and every $0\le k\le g$, together with the
  closed form \eqref{eq:sign}. All values agree.

\item \emph{The character separation.} The grading \eqref{eq:grading} is
  formed for $n\le4$ and the positions of the Weil line and of $\eta^{n}$ are
  computed, confirming \Cref{lem:NSchar}, \Cref{lem:weilchar} and the
  disjointness of \Cref{thm:character}.

\item \emph{The numerical comparison.} The values of $2n^{2}+2n$, $2^{2n}$
  and $3^{2n}$ are tabulated for $2\le n\le8$, confirming
  \Cref{lem:chicompare}; the script also records, as a counterexample to the
  inference that \Cref{rem:retract} retracts, the types $(1,3)$ and $(1,5)$ on
  an abelian surface with their Euler characteristics $3$ and $5$.

\item \emph{The semiregularity target.} The identity
  $\sum_{q}\binom{2n}{q}\binom{2n}{q+2}=\binom{4n}{2n-2}$ of
  \Cref{prop:target} is verified for $2\le n\le7$, with the growth rates of
  \Cref{cor:countnotbind}.

\item \emph{The coordinate model.} The script \texttt{explicit\_weil.py}
  builds the model of \Cref{setup:model} over $\QQ(\sqrt{-d})$ for
  $(n,d)$ in a table covering $n\le3$ and $d\in\{1,2,3,5,7\}$, and checks
  eight things for each: that the $u^{s}_{j}$ of \eqref{eq:eigenvec} are
  eigenvectors of $S$; that $\omega_{+}+\omega_{-}$ and
  $(\omega_{+}-\omega_{-})/\delta$ are rational and agree with the closed form
  \eqref{eq:omegagens}; that $\iota(\tau)^{*}$ multiplies $\omega_{\pm}$ by
  $\tau^{2n}$ and $\bbar{\tau}^{\,2n}$ at five values of $\tau$; that
  $\omega_{\pm}$ is of Hodge type $(n,n)$ for the balanced structure and of
  type $(p,2n-p)$ for every other one, which is \Cref{prop:signatureforced};
  that $\eta$ is rational, of type $(1,1)$ and a norm eigenclass; that the two
  Riemann relations hold with the sign vector of \Cref{setup:model}; that
  $\eta^{2n}\ne0$; and that the supports of $\omega_{1},\omega_{2}$ and
  $\eta^{n}$ are disjoint, which is \Cref{prop:disjoint}.

\item \emph{The Hodge class count.} The script \texttt{hodge\_invariants.py}
  computes, for $n\le4$ and every $(r,s)$, the dimension of the space of
  $\mathfrak{sl}_{2n}$-invariants in
  $\bigwedge^{r}V_{+}\otimes\bigwedge^{s}V_{+}^{\vee}$, as the kernel of the
  raising operators on the weight-zero subspace, in exact rational arithmetic.
  It confirms $\dim\Hdg^{n}=3$ and $\dim\Hdg^{k}=1$ for $k\ne n$, which is
  \Cref{thm:hdgcount}, and the codimension count of \Cref{prop:weildomain}.

\item \emph{The Clifford algebra.} The script \texttt{kugasatake.py} builds
  $C(T,q)$ for forms of signature $(3,b-3)$ with $b\le8$, checks the Clifford
  relations and $v^{2}=q(v)$ for a general $v$, checks $J^{2}=-1$ and that
  left multiplication by $J$ squares to $-\id$ on $C^{+}$, which is
  \Cref{lem:cliffJ}, checks $C^{-}C^{+}C^{-}\subseteq C^{+}$, which is what
  makes $\Psi$ of \Cref{prop:ksembed} land where it should, and tabulates the
  Hodge numbers of \Cref{prop:h2pieces} together with the solution $n=1$ of
  \Cref{prop:normpart}.

\item \emph{The family and the secant plane.} The scripts
  \texttt{weiltype\_family.py} and \texttt{secant\_plane.py} verify
  \Cref{thm:explicitweil} symbolically for $n\le5$ and several $d$, and
  derive the closed forms of \Cref{thm:plane} and
  \Cref{thm:secantinvariants}.

\item \emph{The split member.} The script \texttt{split\_locus.py} builds
  $A=X\times\widehat{X}$ with the action \eqref{eq:Mdef} in coordinates and
  does two independent things. First it solves, by exact rational Gaussian
  elimination, for an expression of each of $\omega_{1},\omega_{2}$ as a
  rational combination of the degree $2n$ monomials in
  $\beta,\widehat{\beta},\ell$, and prints the solution; it succeeds for
  $n\le3$ and $d\in\{1,2,3,7\}$. Second it verifies
  \Cref{thm:splitclosed} directly: that $\gamma\mp\sqrt{-d}\,\ell$ is an
  eigenclass for $\tau^{2}$ respectively $\bbar{\tau}^{\,2}$, and that
  \eqref{eq:splitclosed} holds on the nose, for $n\le5$ and several $d$.

\item \emph{The reach of the base case.} The script
  \texttt{split\_geometry.py} settles the four assertions of
  \Cref{prop:splitgeom}. It verifies $M^{2}=-d$ and $M^{\top}Q+QM=0$ entry by
  entry for the block form $Q=\beta\oplus(\widehat\beta/d)$ on
  $H_{1}(A,\QQ)$, for $n\le4$, and checks that the action induced on
  $H^{2}(A,\QQ)$ multiplies $d\beta+\widehat\beta$ by $d$, multiplies
  $\gamma$ and $\ell$ by $-d$, and has $\beta+d\widehat\beta$ as an eigenclass
  only when $d=1$, for six pairs $(n,d)$; it
  computes $\det H=(-d)^{n}(d_{1}\cdots d_{n})^{2}$ for ten choices of
  $(n,d,\beta)$ including non-principal $\beta$, and certifies in each case
  that $(-1)^{n}\det H$ is a norm from $K$, together with the two facts the
  proof uses, that $\det B$ is a square and that $d$ is a norm; it checks the
  codimension formula $n^{2}-n(n+1)/2=n(n-1)/2$ for $n\le10$; and it computes
  both the dimension of the space of $\mathrm{Sp}$-invariants and the rank of
  the set of monomials $\beta^{i}\widehat{\beta}^{j}\ell^{k}$ in degree $2n$,
  finding rank three in degree two for $n\le4$ and agreement in the middle
  degree for $n\le3$.

\item \emph{Quaternionic multiplication.} The script \texttt{quaternionic.py}
  builds $B=(-d,b)$ and the module $B^{n}$ in coordinates and verifies
  \Cref{thm:quatweil} for nine choices of $(n,d,b,a)$ with $n\le3$: the
  relations $i^{2}=-d$, $j^{2}=b$, $ij=-ji$; that $E$ of \eqref{eq:quatE} is
  alternating and nondegenerate with $i$ anti-self-adjoint and $j$
  self-adjoint; that the hermitian form on the $K$-basis is
  $\mathrm{diag}(2a_{k}d,-2a_{k}bd)$ with signature $(n,n)$; the determinant
  formula and hence the class $b^{n}$; and that $g(x,y)=E(jx,y)$ is
  alternating, $K$-bilinear and nondegenerate, which is the hypothesis of
  \Cref{thm:divcriterion}. It then checks \Cref{thm:everyfamily} numerically:
  for each of several fields, every nontrivial class in the tested range is
  realised in every dimension $n\le8$ with signature $(n,n)$, by a single
  factor when $n$ is odd and by the product of \Cref{lem:products} when $n$ is
  even. Last it records the two dimensions behind \Cref{rem:ntwo}: the
  quaternionic locus has dimension $n(n+1)/2$ inside a period domain of
  dimension $n^{2}$, so its codimension is $n(n-1)/2$, which is checked for
  $n\le10$ and is $1$ exactly at $n=2$.

\item \emph{The semiregularity map of a sum of line bundles.} The script
  \texttt{semiregularity.py} builds the split member with $X=E_{i}^{n}$, so
  that the Hodge decomposition is defined over $\QQ(i)$ and every step is
  exact, and verifies: that the six degree four monomials in
  $\beta,\widehat\beta,\ell$ are independent for $n\le4$, which is
  \Cref{lem:sixmonomials}; that $\theta^{+}\theta^{-}$ is never a multiple
  of $\eta^{2}$, for six pairs $(n,d)$, which is \Cref{lem:tensorrank}; that the coefficient of $\theta^{+}\theta^{-}$
  in $\ch_{2}$ is $\sum_{i}m_{i}\Nm(u_{i})$, which is the identity behind
  \Cref{thm:positivity}; that no multiset with positive multiplicities
  satisfies the resulting equations with some $u_{i}$ nonzero, for five pairs
  $(n,d)$; that the map $\Phi$ of \Cref{lem:splitsemireg}(iii) is injective
  at $n=2$ for $s\le3$ and of rank $22$ rather than $24$ at $s=4$; that
  an exhaustive search over the classes $a\beta+b\widehat\beta+c\ell$ with
  $|a|,|b|,|c|\le3$ and $d\in\{1,2\}$ produces no signed tuple of at most
  three line bundles with a flat Chern character and a nonzero Weil component;
  and that every signed tuple of at most four in the same box with a flat
  Chern character and a nonzero Weil component has exactly four terms and a
  map $\Phi$ with nonzero kernel, of dimension $2$, $3$, $4$ or $6$ at $d=1$
  and $3$ or $6$ at $d=2$. It then turns
  to $n=3$: it finds the Prouhet--Tarry--Escott pair of \Cref{lem:pte} on the
  norm circle $\Nm=4$ of $\QQ(\sqrt{-3})$ by exhaustive search, verifies that
  the six line bundles of \Cref{prop:explicitobject} have
  $\ch_{k}=0$ for $k\ne3$ and $\ch_{3}=12\,\omega_{1}$, and verifies that
  every difference of the six classes is nondegenerate of index three. The
  companion \texttt{semireg\_fast.py} repeats the rank computations modulo a
  prime congruent to $1$ mod $4$, which reaches $n=3$: it reproduces the $n=2$
  ranks exactly, finds $\Phi$ injective at $n=3$ for a general choice of up to
  ten classes, and finds rank $75$ of $90$ for the explicit object.

\item \emph{The annihilator and the rank-zero obstruction.} The script
  \texttt{weil\_annihilator.py} computes, modulo the same prime, the
  dimension of the annihilator of each rational generator of the Weil line
  in $H^{0,2}$ at $n=2,3,4$ and finds $n^{2}$ in every case, which is
  \Cref{lem:annihilator}; splits the semiregularity map of the explicit
  object of \Cref{sec:construction} by parity and finds an odd kernel of
  dimension $9=n^{2}$ equal to the annihilator, an even kernel of dimension
  $6$, and a total of $15$, against $0$ for a generic triple, which is
  \Cref{cor:pteobstructed}; and runs the semiregularity rank over every
  configuration $s_{i}\gamma+c_{i}\ell$ with $\sum s_{i}=\sum c_{i}=0$ in a
  box of side $7$ whose virtual Chern character lies on the Weil line,
  finding eighteen configurations and rank $75$ of $90$ for all of them,
  which is \Cref{rem:familyuniform}.

\item \emph{The Weil class under products.} The script
  \texttt{weil\_product.py} verifies, in exact arithmetic over
  $\QQ(\sqrt{-d})$, the two identities \eqref{eq:weilmult} for every pair
  $n_{1},n_{2}\le3$; their $K$-valued form
  $\omega(A_{1})\omega(A_{2})=2\,\omega(A_{1}\times A_{2})$ over the same
  range; the $k$-factor identity with constant $2^{k-1}$ for six partitions
  with $k=3$ and $k=4$; that the Weil class lies in $H^{2n}$, so that at
  $n=1$ it is a divisor class, which is \Cref{lem:basecase}; and the
  dimension comparison of \Cref{rem:prodscope}, that the product locus has
  dimension $n$ inside a family of dimension $n^{2}$ and is smaller than the
  quaternionic locus for every $n\ge2$.

\item \emph{The annihilator as the tangent space.} The script
  \texttt{weil\_tangent.py} builds a polarised Hodge structure of
  $(K,-1,n)$-Weil type from scratch over the Gaussian rationals, with no
  reduction modulo a prime and no floating point anywhere, and checks five
  things at $n=2,3,4$. That the construction is what it claims to be: the
  polarisation has exactly two nonzero blocks on $V_{+}\times V_{-}$, each a
  perfect pairing of $n$-dimensional spaces, and the two rational generators
  of the Weil line are real of type $(n,n)$. That the annihilator of each
  generator in $H^{0,2}$ has dimension exactly $n^{2}$, which is
  \Cref{lem:annihilator} again, now in a model built independently of the one
  used in item (XIV). That the polarisation-preserving first-order
  deformations have dimension $n(2n+1)=\dim\cA_{2n}$ and those commuting with
  $K$ have dimension $n^{2}=\dim\cD_{n,n}$. That a polarised deformation
  contracts the whole Weil line to zero exactly on that $n^{2}$-dimensional
  subspace, and that one generator already cuts it out, which is
  \Cref{prop:firstorder}. And that the map of \Cref{thm:anntangent} is
  injective, independent of the basis of $V_{+}^{1,0}$ used to write it, and
  has image exactly the annihilator. The five statements are then rerun for
  eight distinct choices of the $n$-plane $V_{+}^{1,0}$ at $n=2$ and at $n=3$,
  and hold in all sixteen.

\item \emph{The quaternionic locus as a Lagrangian Grassmannian.} The script
  \texttt{lagrangian\_\allowbreak locus.py} verifies, over $\QQ$ in exact arithmetic and
  for $2\le n\le5$, the three statements of \Cref{prop:lagrangian}. That the
  form $\omega_{\psi}(x,y)=E(x,\psi y)$ attached to a Rosati-symmetric $\psi$
  is alternating on $V_{+}$, which is the identity $E(\psi x,y)=E(x,\psi y)$
  read twice. That the graph of a symmetric form is isotropic and the graph of
  a form that is not symmetric is not, so the chart of
  $\mathrm{LG}(n,2n)$ by symmetric forms is exactly right and not merely
  contained in the locus. That the chart has dimension $n(n+1)/2$, computed as
  the rank of the symmetric basis rather than quoted, and that the codimension
  in $\cD_{n,n}$ is therefore $n(n-1)/2$; at $n=4$ these are $10$ and $6$. And
  that a random $n$-plane, completed to a basis of $V_{+}$, is Lagrangian for
  the form that is standard in that basis, which is the construction behind
  \Cref{rem:onlyrational}: over the reals the loci sweep the whole family, and
  only the rationality of $\psi$ prevents them from covering it.

\item \emph{The size forced on an object carrying the class.} The script
  \texttt{object\_\allowbreak size.py} verifies, in the exterior algebra over $\QQ$ for
  $1\le n\le4$, that $\alpha_{\pm}^{2}=0$ and hence
  $\omega_{1}^{2}=2\alpha_{+}\alpha_{-}$, that $\int\omega_{1}^{2}\ne0$, and
  that $\chi(E,E)=(-1)^{n}N^{2}\int\omega^{2}$ for every rank $r$ and every
  multiplicity $N$ in the ranges tested, the rank cancelling out of the
  identity. That is \Cref{prop:objectsize}, and with it the lower bound
  $\sum_{i}\dim\Ext^{i}(E,E)\ge N^{2}|\int\omega^{2}|$. The script also records
  that the constant is not dimension-free, since it depends on the
  normalisation of $\omega$ against $\eta$, so that no growth rate in $n$
  should be read off it; what the bound needs, and all that it needs, is that
  the Weil class is not isotropic.

\item \emph{Infinitesimal rigidity of divisor classes.} The script
  \texttt{rigidity.py} builds the four blocks of $H^{1,1}$ and the tangent
  space of $\cD_{n,n}$ over $\QQ(i)$ and verifies, for $n=2,3,4$, the three
  parts of \Cref{prop:rigidity}. That the real span of the classes killed by
  contraction has dimension exactly one and is the polarisation, checked
  against the full real space of $(1,1)$ classes, of dimension $4n^{2}$,
  which is $16$ at $n=2$ and $64$ at $n=4$. That the kernel attached to a
  norm-character class with block $a$ has dimension $n(n-\operatorname{rank}a)$
  for every rank; at $n=4$ the four values are $12$, $8$, $4$, $0$. And that a
  class drawn at random has kernel zero, so it rigidifies the family to a
  point. The dimension $n^{2}$ of the tangent space is recovered in the same
  computation rather than assumed.

\item \emph{The invariants of the required object.} The script
  \texttt{integrality.py} computes, for $1\le n\le8$, every squarefree
  $d\le11$ and every rank and twist in a range, the Chern classes of the
  Chern character $au+bv$ by Newton's identities, and finds every one of them
  integral; it also confirms the closed form
  $\Delta=a^{2}d+b^{2}$ of \Cref{thm:secantinvariants} and its positivity.
  That is \Cref{prop:notforbidden}: nothing numerical forbids the object the
  construction needs, in any dimension tested.

\item \emph{A sheaf with the required Chern character exists.} The script
  \texttt{secant\_\allowbreak exists.py} carries out, exactly over $\QQ$ and for
  $1\le n\le10$, the computation behind \Cref{thm:secantexists}. For every
  squarefree $d\le11$ and every rank and twist in a range it solves the
  Vandermonde system that expresses the rational point $au_{t}+bv_{t}$ of the
  secant plane as a rational combination of the Chern characters
  $e^{jt}=\ch(\cO_{X}(jt))$, clears denominators to obtain the least positive
  multiple $M$ and the integer witness \eqref{eq:lbwitness}, and then verifies
  that witness by recomputing every moment $\sum_{j}n_{j}j^{k}$ and comparing
  it with $Mc_{k}$, rather than trusting the solve. It checks that the rank
  $\sum_{j}n_{j}$ of the witness equals $Ma$, which is the positive rank
  hypothesis of \Cref{lem:posrank}, and it records the least $M$ in each case:
  $6$ at $n=4$ with $d=3$ and $a=b=1$, at most $2016$ over the whole range at
  $n=8$, and never $1$ once $n\ge7$, which is \Cref{rem:themultiple}. Finally
  it expands $\theta^{k}/k!$ in the exterior algebra of a symplectic lattice
  and finds every coefficient integral, which is why the coordinates $c_{k}$
  are integers and hence why the problem is a lattice problem at all.

\item \emph{Which split objects can be semiregular.} The script
  \texttt{pte\_\allowbreak search.py} settles \Cref{thm:pterange} in exact integer
  arithmetic. It enumerates the norms represented by $\cO_{K}$ below a bound,
  forms for every set of $n+1$ of them the forced level sums
  \eqref{eq:levelsums} and the least integral scale, and tests them against
  the supply $r_{K}(N_{j})$ of elements of that norm; at $n+2$ norms the
  solution space is two dimensional and its integer points in the same box are
  swept coordinate by coordinate. For $4\le n\le8$ and $d\le19$ nothing
  survives, with norms to $110$ at $n=4$ and to $30$ at $n=8$. At $n=3$ and
  norms to $200$ exactly eleven sets survive, all at $d=1$ or $d=2$, of the
  three shapes recorded in \Cref{thm:pterange}(iii). The script then runs the
  whole system \eqref{eq:pteconditions} on those sets: for each it enumerates
  every assignment of signs to the elements of $\cO_{K}$ at those norms with
  the prescribed level sums, splitting the four levels into a tabulated half
  and a streamed half so that the enumeration is complete, and finds no
  solution for eight of them. The companion script
  \texttt{pte\_\allowbreak remaining.py} settles the other three, $d=1$ at
  $(13,52,117,130)$, $(17,68,153,170)$ and $(25,50,100,125)$, whose largest
  level has sixteen elements: it streams three levels against the fourth
  through a linear hash of the ten off-diagonal coordinates, re-checks every
  match exactly, and finds no solution. It also reruns two of the eight as a
  control. The largest search covers $4.29\times10^{14}$ sign patterns.
\item \emph{The quaternionic Weil cycle.} The script
  \texttt{quaternionic\_\allowbreak divisibility.py} checks
  \Cref{prop:quatdivisible} on the lattice $\cO_{b}^{2}$. It confirms the
  ratios $-b/3$ and $4b^{2}/3$ and the value $1/12$ for $d\in\{1,3\}$, two
  weight vectors and fifteen values of $b$ up to $101$; the decomposition
  $E_{b}=E_{u}+bE_{v}$, $g_{b}=bg_{1}$, $\lambda_{b}=b\lambda_{1}$ with $g_{1}$
  and $\lambda_{1}$ integral; the equalities $Z_{b}=b^{2}Z_{1}$ and
  $Z'_{b}=b^{2}Z'_{1}$ in $\wedge^{4}\Lambda^{*}$; the growth
  $\int\eta_{b}^{4}=b^{2}\int\eta_{1}^{4}$; and the effect of replacing $j$ by
  $\alpha j$. It then computes the saturation of the Weil plane in
  $\wedge^{4}\Lambda^{*}$ and the sublattice spanned by products of two
  integral divisor classes, for $d\in\{1,2,3,7\}$, three weight vectors and
  seven values of $b$, and finds Gram matrix $\mathrm{diag}(8d,8d^{2})$ and
  index $2(a_{1}a_{2})^{2}$ throughout.
\item \emph{The divisor route on a fixed lattice.} The script
  \texttt{divisor\_\allowbreak route.py} checks \Cref{prop:pencil} on
  $\Lambda=\cO^{2}$ with $j^{2}=1$, for $d\in\{1,3\}$. It computes the
  lattice of $K$-bilinear classes, of rank $12$, checks $t(x)=0$ on it and the
  identity $-\int xy\,\eta^{2}\big/\!\int\eta^{4}=\mathrm{tr}(\phi_{x}\phi_{y})/24$
  on all $78$ pairs of basis vectors, and finds signature $(8,4)$ for $I$. It
  verifies that $\phi_{x}^{2}$, $\phi_{y}^{2}$ and
  $\phi_{x}\phi_{y}+\phi_{y}\phi_{x}$ are scalar for the pencil, and
  rechecks $\phi_{x+ky}^{2}=\beta(k)$ and $I(x+ky)=\beta(k)\int\eta^{4}/3$
  directly at eight values of $k$. For $0\le k\le10$ it saturates
  $\QQ\eta+\QQ x_{k}+\QQ x_{k}(i\cdot,\cdot)$ in $H^{2}(\Lambda,\ZZ)$,
  checks that $I$ is positive definite on the result, and computes the
  minimum $\mu_{k}$ of $I$ off $\QQ\eta$ by an exhaustive search in a box
  derived from the inverse Gram matrix. For the bound valid at every $k$ it
  forms the $2\times2$ minors of the matrix with rows $x_{k}$ and
  $x_{k}(i\cdot,\cdot)$ as polynomials in $k$, computes resultants of pairs of
  them, and obtains $N_{1}=2$ and $N_{3}=4$; it checks directly that the gcd of
  the minors divides $N_{d}$ for $|k|\le200$. Finally, for $0\le k\le4$ it
  computes the centraliser of $i$ and $\phi_{x_{k}}$ in $\mathfrak{sp}(V,E)$,
  of dimension $10$, and its invariants in $\bigwedge^{2}V^{*}$, of dimension
  $3$; and at the one value of $k$ in $[-6,6]$ for each $d$ at which
  $\beta(k)$ is a rational square it constructs a rational complex structure
  $J_{0}$ on the pencil member, commuting with $i$ and $\phi_{x_{k}}$ and
  compatible with $E$, and checks that $[\mathfrak{c},J_{0}]$ has dimension
  $6$ and that $[\mathfrak{c},J_{0}]$ together with its brackets spans all of
  $\mathfrak{c}$, which is the Lie algebra step in the proof of
  \Cref{prop:pencil}(iv).
\item \emph{The obstruction map of a split object.} The script
  \texttt{split\_\allowbreak obstruction.py} works in the model of item
  (XIII), modulo the same prime, and checks \Cref{prop:splitobstruction}.
  For $(n,d)=(2,1),(2,3),(3,3)$ it computes the space $T$ of first-order
  deformations killing $\eta$ and both generators of the Weil line and
  finds dimension $n^{2}$; computes the space $T_{\mathrm{spl}}$ of
  deformations keeping $\beta,\widehat\beta,\ell$ of type $(1,1)$ and finds
  dimension $n(n+1)/2$ and $T_{\mathrm{spl}}\subseteq T$; as a control,
  replaces $\eta$ by $\beta+d\widehat\beta$ and finds that $T$ collapses to
  dimension $n(n+1)/2$ for $d>1$, which is the mismatch of conventions
  described in the footnote to \Cref{setup:construction}; and for three
  classes spanning $\NS(A_{0})_{\QQ}$ together with $\eta$ computes the
  obstruction map $v\mapsto(v\lrcorner\lambda_{i})_{i}$ on $T$, finding
  kernel exactly $T_{\mathrm{spl}}$ and rank $n(n-1)/2$; it also finds that
  each of the seven single classes $\gamma$, $\ell$, $\gamma+\ell$,
  $\gamma+2\ell$, $\beta$, $\widehat\beta$, $\beta+\ell$ already has
  $\{v\in T:v\lrcorner\theta=0\}=T_{\mathrm{spl}}$, and that each of $150$
  random rational classes off $\QQ\eta$ is contracted nontrivially by some
  vector of $T$ while $\eta$ and $3\eta$ are contracted to zero by all of
  $T$, which is the mechanism of \Cref{thm:linebundle}. For the six line
  bundles of \Cref{prop:explicitobject} at $n=3$ it finds kernel
  $T_{\mathrm{spl}}$ of dimension $6$ and rank $3$; checks that
  $\sigma_{E}$ kills the image; computes the annihilator of $\omega_{1}$ in
  $H^{0,2}$, of dimension $9$, and checks that $\sigma_{E}$ kills
  $z\cdot\id_{E}$ for $z$ in it; finds that the sum of that
  nine-dimensional space and the three-dimensional image has dimension
  $12$; recomputes the kernel of $\sigma_{E}$ as $15$; and checks that the
  components of every obstruction vector on $L_{i}$ and on $L_{i}^{-1}[3]$
  are negatives of each other.
\item \emph{The evaluation map.} The script
  \texttt{evaluation\_\allowbreak map.py} rewrites the model of item (XIII)
  in its Hodge basis and checks \Cref{prop:evaluation} for the explicit
  object at $n=3$. It builds a basis of $HT^{2}$, of dimension
  $15+36+15=66$, implements the three actions of \eqref{eq:HT2} as
  multiplication, a derivation and a double interior product, and verifies
  the identity $\sigma_{E}(\operatorname{ev}_{E}(x))=x\lrcorner\ch(E)$ on
  every basis vector and on random elements; confirms that $\ch(E)$ is
  concentrated in degree six and proportional to $\omega_{1}$; computes the
  rank of $\operatorname{ev}_{E}$, which is $45$, against
  $\dim\Ext^{2}(E,E)=90$; computes the three kernels of contraction with
  $\ch(E)$, of dimensions $9$, $18$ and $9$; finds that their images under
  $\operatorname{ev}_{E}$ have dimensions $9$, $6$ and $9$, that the first and
  third coincide, that the first and second meet in zero, and that the sum of
  all three has dimension $15$, equal to the dimension of $\ker\sigma_{E}$,
  which it recomputes; and repeats the last comparison for every one of the
  eighteen configurations of item (XIV), finding $\dim\ker\sigma_{E}=15$ and
  $\dim(\ker\sigma_{E}\cap\image\operatorname{ev}_{E})=15$ in every case.
\item \emph{Markman's candidate in dimension eight.} The script
  \texttt{markman\_\allowbreak candidate.py} works in
  $\QQ[\Theta]/(\Theta^{5})$ with $\int\Theta^{4}=24$ and checks
  \Cref{prop:markman8} for every odd squarefree $d\le101$: the Euler
  characteristics $14$, $-36$ and $24$ from the Koszul complexes; the count
  $12N(N-5)=(d+9)(3d-3)$ of isolated points and its comparison with the
  $(d+9)(2d-1)$ normalised; $\chi(\cO_{\widetilde Z})=(d+9)(27-d)$ from the
  Mayer--Vietoris complex; the equality $\ch(\cF)=u_{t}+3v_{t}$ degree by
  degree for the twist $3\Theta$, and the failure by $-4\Theta^{2}$ in
  degree two for the twist $\Theta$; $\Nm(3+\sqrt{-d})=2N$ and the
  discriminant $(d+9)\Theta^{2}$; and the ranks $8d(d-9)$ and $8d(d+9)$ of
  the two transforms, with the threefold control $-8q$ that reproduces
  Markman's rank at $n=3$.
\item \emph{The classes preserving a secant Chern character.} The script
  \texttt{secant\_\allowbreak kernel.py} works in the Hodge algebra of an
  abelian $n$-fold with $t=\sum_{a}p_{a}\wedge q_{a}$, modulo a large prime,
  and checks \Cref{thm:factor}(i), (ii) and \Cref{lem:bfield}(i) for
  $(n,d,a,b)$ equal to $(3,3,1,1)$, $(3,3,1,3)$, $(4,3,1,3)$, $(4,5,1,3)$,
  $(4,7,2,1)$ and $(5,3,1,3)$, the third being Markman's candidate. It builds
  $\ch(F)=au_{t}+bv_{t}$, implements the actions of $HT^{1}$ and $HT^{2}$ as
  multiplication, interior product, derivation and double interior product,
  and finds: that the kernel of $y\mapsto y\lrcorner\ch(F)$ on $HT^{1}$, of
  dimension $2n$, is zero in every case; that the kernel on $HT^{2}$ has
  dimension exactly $n^{2}$, contains the $n(n+1)/2$ polarised deformations
  and the $n(n-1)/2$ classes $(\tfrac d2\,\pi\lrcorner t^{2},0,\pi)$, and is
  spanned by them; and, as controls, that the bare classes $(0,0,\pi)$ and
  the classes $(-\tfrac d2\,\pi\lrcorner t^{2},0,\pi)$ with the compensation
  of the wrong sign all act nontrivially on $\ch(F)$. It then verifies the
  identity $x\lrcorner(we^{B})=((e^{B}x)\lrcorner w)e^{B}$ of
  \eqref{eq:bfield} on random even classes $w$, random $x\in HT^{2}$ and
  $B=kt$ for $n=3,4$, and the count $n(n+1)/2+n(n-1)/2=n^{2}$ for $n<40$.
\item \emph{The local products at a double point.} The Macaulay2 script
  \texttt{m2/local\_\allowbreak products.m2} recomputes
  \Cref{lem:localproducts} over $\QQ[x_{1},\dots,x_{4}]$. For each of the
  two local models, the ideal $I=(x_{1},x_{2})\cap(x_{3},x_{4})$ and the
  complex $K=[R\to R/(x_{1},x_{2})\oplus R/(x_{3},x_{4})]$, resolved by the
  cone of the lifted map, it builds the jet chain maps $-\partial_{a}(d)$ on
  the free resolution, checks that they anticommute with the differential
  and are not null homotopic, and for every pair $(a,b)$ decides whether the
  composite is null homotopic. It finds $\Ext^{2}=k^{4}$ for $I$, with the
  four mixed products nonzero and spanning it and the same-plane products
  zero; $\Ext^{2}=k^{2}$ for $K$, with the two same-plane products nonzero
  and spanning it and the mixed products zero; anticommutation
  $a_{a}a_{b}+a_{b}a_{a}=0$ in both models; and, for $I$, the same vanishing
  pattern from the jet modules and \texttt{yonedaProduct}. This item needs
  Macaulay2 (version 1.22 with the package \texttt{Complexes}) and is not
  part of \texttt{verify\_all.py}; its transcript
  \texttt{m2/local\_products.txt} is included with the code.

\item \emph{A rank one secant object with smooth support.} The script
  \texttt{smooth\_\allowbreak support.py} expands $\ch(\cO_{S})=1-\ch(\cF)e^{-b\Theta}$
  from scratch in $\QQ[\Theta]/(\Theta^{5})$, for every $b\le6$ and every
  squarefree $d<60$, and reads off the class $N\Theta^{2}$ with
  $N=\Nm_{K/\QQ}(b+\sqrt{-d})/2$, the pushforward
  $\iota_{*}K_{S}=(4bN/3)\Theta^{3}$ and the Euler characteristic
  $\chi(\cO_{S})=4N(2b^{2}-N)$, checking each against the closed form rather
  than quoting it. It then verifies that Noether's formula and the
  self-intersection formula together give $K_{S}^{2}=12N(4b^{2}-N)$ and
  $e(S)=12N(4b^{2}-3N)$, that the Hodge index bound $9N\ge4b^{2}$ is implied
  by $d>0$ and so never binds, that $e(S)\ge0$ is $d\le\tfrac53b^{2}$ and that
  $3e(S)-K_{S}^{2}=24(b^{4}-d^{2})\ge0$ is $d\le b^{2}$, and that equality in
  the Hodge index would need $9N=4b^{2}$, which no admissible pair reaches.
  At $b=3$ it prints the six surviving discriminants, the four of them with
  $N$ an integer, and \Cref{tab:smoothsupport}. That is
  \Cref{thm:smoothinvariants} and \Cref{cor:fourdiscriminants}.

\item \emph{Where the local obstruction lives.} The Macaulay2 script
  \texttt{m2/lci\_\allowbreak products.m2} computes, for ten codimension two germs in
  $\CC^{4}$, the number of minimal generators, the projective dimension of
  $R/I$, the ranks of a minimal free resolution, and
  $\cE xt^{1}$, $\cE xt^{2}$ and $\cE xt^{3}$ of the ideal against itself; for
  the germs with $\cE xt^{2}\ne0$ it also decides, by null homotopy of the
  composite of the two jet chain maps, which products of translation classes
  are nonzero. The seven Cohen--Macaulay germs, namely the smooth one, two
  and three planes meeting along a line, the double plane, the quadric cone,
  the fat plane and the cone over the twisted cubic, all have
  $\cE xt^{2}=0$, the last two despite needing three generators in codimension
  two. The three that are not Cohen--Macaulay, namely two planes meeting at a
  point, three planes meeting pairwise at a point and a plane with an embedded
  point, have $\cE xt^{2}\ne0$ and nonzero products. That is
  \Cref{thm:cmvanishing} and \Cref{rem:cmsharp}. The script runs separately
  from \texttt{verify\_all.py} because it needs Macaulay2; its transcript is
  distributed with it.

\item \emph{Cohen--Macaulay supports with a split resolution.} The script
  \texttt{hilbert\_\allowbreak burch.py} reads the coefficients
  $(1,b,-d,-db,d^{2})$ off $u_{\Theta}+bv_{\Theta}$ rather than quoting them,
  confirms the three identities $c_{k+2}+dc_{k}=0$ that make the weighted
  measures of \Cref{thm:noci} share mass, mean and variance, and checks that
  the quadratic of that theorem has discriminant $-\tfrac29b^{2}-2d<0$ for
  every $b\le11$ and every squarefree $d<200$, so that no complete
  intersection can occur. It then enumerates \eqref{eq:burchsystem} exactly,
  by meet in the middle on the four moments, over every multiset of twists in
  the boxes of \Cref{tab:burch}, imposing the two conditions of
  \Cref{lem:burchmoments}. Nothing survives at $r=2$ or $r=3$, for $b=3$ and
  every squarefree $d<24$ and also for every $b\le6$ over seven fields; at
  $r=4,5,6$ only $d=15$ and $d=23$ survive, and the first candidate at each
  rank is printed. The script also records that the moment system alone is not
  enough: at $r=4$, $d=23$ it has the solution $p=(-8,-5,0,1,2)$,
  $q=(-7,-6,-4,4)$, which only the injectivity of $\varphi$ removes.

\item \emph{The closure graph.} The script \texttt{closure\_\allowbreak graph.py}
  carries the rule set of \Cref{setup:rules} as data: thirty-seven statements,
  each marked proved here, quoted, or open, and seventeen rules, each carrying
  the label of the theorem that proves it. It checks that every statement
  named in a rule is declared and every declared statement is used, that the
  rule set is acyclic, and that every one of the fifteen theorem labels it
  names, for its rules and for the consequences of $\mathbf{HC}$ it records,
  is a label of this paper, against the list of all $534$ labels generated
  from the sources and distributed with the script as
  \texttt{paper\_\allowbreak labels.txt}. It then computes the forward closure
  by iterating the one step map to a fixed point, and verifies that no
  statement marked proved rests at any depth on an open one, that
  $\mathbf{HC}$ is not in the closure of the proved and quoted statements,
  and, by exhaustive search over all $2^{13}$ subsets of the thirteen open
  statements, that there are exactly the five minimal sufficient sets of
  \Cref{thm:closure}(ii), $\{\mathrm{F3}\}$, $\{\mathrm{L},\mathrm{M}\}$,
  $\{\mathrm{L},\mathrm{F3}'\}$, $\{\mathrm{V},\mathrm{F3}'\}$ and
  $\{\mathrm{P2},\mathrm{F2},\mathrm{F3}'\}$, that (F3) is the one open
  statement that suffices alone, that \textup{(F3$'$)} alone gives neither
  the conjecture for abelian varieties nor $\mathbf{HC}$, that the sets
  $\{\mathrm{P2},\mathrm{F2},\mathrm{F3}\}$, $\{\mathrm{L},\mathrm{F3}\}$
  and $\{\mathrm{V},\mathrm{F3}\}$ suffice and are not minimal, that the
  route through this paper is the only minimal set with three elements and
  the only one whose derivation uses a statement proved here, that each
  element of each set is necessary, that every set contains (F3),
  \textup{(F3$'$)} or \textup{(M)}, that every element of every set other
  than (P2) is recorded as a consequence of $\mathbf{HC}$, and that
  $\{\mathrm{F3}\}$ and $\{\mathrm{L},\mathrm{M}\}$ are each equivalent
  to it. It confirms that \textup{(V)} alone gives the
  conjecture for abelian varieties and not $\mathbf{HC}$, and that
  \textup{(L)} gives \textup{(V)}. It confirms that the six open demands of
  the secant route lie outside every minimal set, that granting the two of
  them that matter reaches $\mathbf{W}(K,n,\delta_{0})$ and not the rest, that
  the propagation statement reaches all of it and also the Weil classes of
  every CM field, so that those are not open, and that the bounded criterion
  is absent from the rule set, being equivalent to its own conclusion. It
  then removes the rule of \Cref{prop:f3ishc} and the statement
  \textup{(F3$'$)} with its rule and repeats the search, which returns the
  four sets of the earlier version of this paper,
  $\{\mathrm{L},\mathrm{M}\}$, $\{\mathrm{L},\mathrm{F3}\}$,
  $\{\mathrm{V},\mathrm{F3}\}$ and $\{\mathrm{P2},\mathrm{F2},\mathrm{F3}\}$;
  so the change in the answer is those two additions and nothing else.
  Finally it checks that the numbering under which the same graph is carried
  into Section 24 of the Lean certificate maps these rules onto those. The
  full derivation of $\mathbf{HC}$ from the proved and quoted statements and
  each minimal sufficient set is printed, one rule per line. That is
  \Cref{thm:closure}, \Cref{prop:f3ishc}, \Cref{prop:f3prime} and
  \Cref{cor:fourremain}.
\item \emph{No split resolution runs the criterion.} The script
  \texttt{split\_\allowbreak resolution.py} records the data behind
  \Cref{thm:nosplit}: that the scalar $(e^{2}+d)/2$ by which the compensated
  Poisson class acts on $\cO(e\Theta)$ is nonzero for every twist and every
  $d>0$ in a range; the vanishing pattern $H^{2}(\cO(m\Theta))=0$ for $m\ne0$
  on an $n$-fold with $3\le n\le6$, from the index theorem; that the three
  split resolutions of \Cref{tab:burch} have no twist in common between
  $E_{0}$ and $E_{1}$ and satisfy the moment identities \eqref{eq:burchsystem},
  so that the theorem removes them as they stand; and, as a control, that no
  Koszul resolution of a complete intersection, twisted by any $b\in[-6,12]$,
  is secant for $d<30$, which is \Cref{thm:noci} in that range.
\item \emph{The two statements that are not reductions, and what the rank
  forces.} The script \texttt{exceptional\_\allowbreak classes.py} has three
  parts. For an abelian fourfold $X$ of Mumford type it decomposes every
  $\bigwedge^{k}V$, $V=V_{1}\otimes V_{2}\otimes V_{3}$, under
  $\mathrm{SL}_{2}^{3}$ by peeling highest weights, finds one invariant in
  $\bigwedge^{2}V$ and one in $\bigwedge^{4}V$, counts eight invariants in
  $\bigwedge^{4}(V\oplus V)$ distributed $1,1,4,1,1$ over the K\"unneth
  summands and three in $\bigwedge^{2}(V\oplus V)$, multiplies the three
  divisor classes $\psi_{11},\psi_{12},\psi_{22}$ in the exterior algebra of
  $V\oplus V$ and finds their six products independent, and counts the
  $\Sp_{8}$-invariants of $\bigwedge^{4}(V\oplus V)$ by the Weyl character
  formula over the $384$ signed permutations, finding six. That is
  \Cref{prop:mumfordproduct}. For the quintic threefold it computes the
  Hodge numbers $(1,101,101,1)$ from the Jacobian ring of the Fermat quintic,
  checks in a four-dimensional model that the map from the $(0,3)$ part to
  the $(3,0)$ part is symplectic of adjoint weight $3$, and lists the adjoint
  weights, $\{-1,0,1\}$ in weight one against $\{-3,\dots,3\}$ here. That
  is \Cref{prop:noabeliantype}. For $n=2$ and $n=3$ it builds
  $H^{*}(A,\CC)$ with the eigenspace decomposition, the Weil class
  $\alpha_{+}+\alpha_{-}$, and the action of all of $HT^{2}$ on it, finds the
  annihilator of dimension $4n^{2}$ spanned by the mixed block, the $K$-linear
  deformations and the mixed bivectors with the remaining classes acting
  injectively, computes the ideal generated by the mixed block in the exterior
  algebra on $2n$ generators, of codimension $2^{n+1}-1$ with quotient of
  dimension $2\binom{n}{k}$ in degree $k\ge1$, and checks over $\mathbb{F}_{3}$ in
  dimensions two and three that two subspaces spanning the space with all
  mixed wedge products zero cannot both be nonzero. That is
  \Cref{thm:p2forces}.
\item \emph{The Weil classes of a CM field of degree four and six.} The
  script \texttt{cm\_\allowbreak fields.py} builds the cyclotomic fields
  $\QQ(\zeta_{5})$, $\QQ(\zeta_{7})$, $\QQ(\zeta_{8})$ and $\QQ(\zeta_{9})$
  exactly, as $\QQ[x]/(p)$ with elements stored as rational coordinate
  vectors, and carries out the construction of \Cref{thm:cmbasepoint} for six
  pairs $(F,n)$, with $m=2$ and $m=3$ and $n=1,2,3$. For each it checks that
  the CM type and its conjugate partition the embeddings, that
  $\zeta-\zeta^{-1}$ is totally imaginary, that the trace-dual bases satisfy
  $\sum_{k}\sigma(\omega_{k})\sigma'(\omega^{*}_{k})=\delta_{\sigma\sigma'}$
  over every pair of embeddings, that $V^{1,0}$ is half-dimensional and
  isotropic for $E=\Tr(\xi H)$ so that the point lies in the period domain,
  that every eigenspace meets $V^{1,0}$ in dimension $n$ so that the Weil
  classes are of type $(n,n)$, that
  $\delta_{i}(f)=\sum_{\sigma}\sigma(f)\,u_{i,\sigma}\wedge u_{i',\sigma}$
  and is therefore of type $(1,1)$, that every $\alpha_{\sigma}$ is nonzero,
  and that the balanced $n$-fold product of the $\delta_{i}$ is
  $w(f)=\sum_{\sigma}\sigma(f)\alpha_{\sigma}$, rational, spanning $W_{F}$.
  For $F=\QQ(\zeta_{5})$ at $n=2$ it also computes, in the exterior algebra of
  the sixteen-dimensional space, the space of rational $(1,1)$ classes, of
  dimension $32$, and records that the balanced sum is necessary: the plain
  product $\delta_{1}(1)\delta_{2}(1)$ is not a Weil class. Finally, for
  $F=\QQ(\zeta_{8})\supset K=\QQ(\sqrt{-1})$ it checks that each $K$-Weil class
  is the product of the two $F$-Weil classes above it and that $W_{K}$ lies in
  the span of the products of pairs from $W_{F}$, which is
  \Cref{prop:composite}.
\item \emph{What the transport does along the orbit.} The script
  \texttt{transport\_\allowbreak growth.py} works on the lattice
  $\Lambda=\ZZ^{2G}$ with the standard alternating form. For a sample of
  rational symplectic elements drawn from the unipotent radical of the
  stabiliser of an isotropic line, with denominators up to $29$, it verifies
  the three scaling identities of \Cref{lem:transport}, namely
  $\det\varphi=c^{2G}$, $\varphi^{*}E=c^{2}E$ and
  $\bigwedge^{k}\varphi=c^{k}\bigwedge^{k}g$, the last on the exterior
  algebra for $G=2,3,4$, and records that the denominators are unbounded. It
  then takes two sublattices of rank $2n=4$ in an abelian fourfold, one
  preserved by the sample elements and one not, computes the multiplicity
  $m=[\varphi^{-1}(\Lambda)\cap W:\Lambda_{W}]$ by Smith normal form and the
  degrees as Pfaffians, and checks
  $\deg\varphi(B)=c^{2n}\deg(B)/m$ in every case. The two cases separate: for
  the preserved subtorus $m=c^{2n}$ exactly and the image degree is constant,
  while for the other $m=c^{2n-1}$ and the image degree grows linearly in $c$.
  Finally it tabulates the bound $c^{2n}\deg(Z)/|G(Z)|$ of
  \Cref{thm:p1forces}(ii) for stabilisers of order $1$, $2$ and $4$, which is
  strictly increasing and unbounded. That is \Cref{lem:transport},
  \Cref{lem:multiplicity} and \Cref{thm:p1forces}.
\item \emph{The annihilator of a Weil class in Hochschild cohomology.} The
  script \texttt{hochschild\_\allowbreak annihilator.py} builds
  $HH^{*}(A)=\bigwedge^{*}\bigl(H^{0,1}\oplus H^{0}(T_{A})\bigr)$, of
  dimension $2^{4n}$, as an exterior algebra on $4n$ operators indexed by bit
  masks, acting on $H^{*}(A,\CC)$ by wedging and contracting, and computes the
  annihilator of $\omega=a\alpha_{+}+b\alpha_{-}$ in every degree by exact
  rational elimination. It confirms that $P=V_{+}^{0,1}\oplus T_{-}$ is
  exactly the annihilator of $\alpha_{+}$ in degree one and
  $Q=V_{-}^{0,1}\oplus T_{+}$ that of $\alpha_{-}$, that
  $\Ann_{HH^{1}}(\omega)=0$, that $\Ann_{HH^{2}}(\omega)=P\wedge Q$ of
  dimension $4n^{2}$, that
  $\dim\Ann_{HH^{k}}(\omega)=\binom{4n}{k}-2\binom{2n}{k}+\delta_{k,2n}$
  in every degree, that the ideal generated by $P\wedge Q$ consists of the
  monomials meeting both halves and has codimension $2^{2n+1}-1$, that
  contraction with $\alpha_{-}$ is injective on $\bigwedge^{*}P$ and with
  $\alpha_{+}$ on $\bigwedge^{*}Q$ with the two images meeting in exactly one
  line in degree $2n$, and that
  $\dim HH^{2}-\dim\Ann_{HH^{2}}=2n(2n-1)$. The computation is run in every
  degree for $n=1$, where degree two is the exceptional degree $2n$ and the
  annihilator is one larger, and for $n=2$ with two choices of $(a,b)$, and in
  degrees up to three for $n=3$. That is \Cref{thm:hhsplitting} and
  \Cref{cor:hhfactor}.

\item \emph{The criterion as a number, and the support of an object meeting
  it.} The script \texttt{p2\_\allowbreak support.py} checks the finite
  linear algebra of \Cref{thm:p2numerical}, \Cref{lem:finiteproduct} and
  \Cref{thm:p2support} in five parts. (A) In the exterior algebra of
  $H^{*}(A,\CC)$ with $HH^{1}$ acting by wedging and contracting, the map
  $\xi\mapsto\xi\lrcorner\,\omega$ on $HH^{2}$ has rank exactly $2n(2n-1)$
  for $n=2$ and $n=3$, that is $12$ and $30$, and it is injective on
  $\bigwedge^{2}P\oplus\bigwedge^{2}Q$; and
  $\binom{4n}{2}-4n^{2}=2\binom{2n}{2}=2n(2n-1)$ for every $n\le60$. (B) On a
  torus of dimension $c$, for $2\le c\le6$, the class of a point contracted
  with $a\wedge b$ is nonzero for all coordinate pairs and for twenty-five
  random rational pairs of independent vectors, and zero for twenty-five
  dependent pairs. (C) For $n=2,3,4$ and random subspaces $W$ of
  $T_{+}\oplus T_{-}$ of codimension at least two, together with the boundary
  cases where $W$ contains $T_{+}$ or $T_{-}$, the images of $T_{+}$ and
  $T_{-}$ in the quotient have all mixed wedges zero exactly when $W$
  contains $T_{+}$ or $T_{-}$; this is the dichotomy used in the proof of
  \Cref{thm:p2support}(i). (D) In an explicit rational model of $\RR^{4n}$
  with complex structure $J$ and an action $M$ of $\QQ(i)$ of signature
  $(n,n)$, scrambled by a random rational change of basis, the forms
  vanishing on $\ker(J-M)$, the real form of $T_{+}$, are exactly
  $V_{-}^{1,0}\oplus V_{+}^{0,1}$, computed over the Gaussian rationals for
  $n=1,2$. (E) For $2\le n\le6$ neither $\alpha_{+}$ nor $\alpha_{-}$ is a
  monomial of $\bigwedge^{*}(V_{-}^{1,0}\oplus V_{+}^{0,1})$ or of
  $\bigwedge^{*}(V_{+}^{1,0}\oplus V_{-}^{0,1})$, so the Weil line meets the
  sum of the two subrings in zero.

\item \emph{Two jet classes on a complex of finite length.} The Macaulay2
  script \texttt{m2/finite\_\allowbreak length\_\allowbreak products.m2}
  builds, for a bounded complex $M$ of free modules over $\QQ[x_{1},x_{2}]$ or
  $\QQ[x_{1},x_{2},x_{3}]$ with finite length cohomology, the jet chain maps
  $-\sum_{i}a_{i}\partial_{i}(d)$ and decides whether the composite of two of
  them is null homotopic. It first confirms that each input has finite
  length, and then tests every pair of coordinate directions and a random
  pair of independent rational directions on eighteen modules: in two
  variables the residue field, $R/\mathfrak{m}^{2}$, $R/\mathfrak{m}^{3}$,
  four monomial ideals, the complete intersection
  $(x_{1}^{2}+x_{2}^{2},x_{1}x_{2})$, the ideal
  $(x_{1}^{3}-x_{2}^{3},x_{1}^{2}x_{2},x_{1}x_{2}^{2})$, the non-cyclic
  modules $\mathfrak{m}/\mathfrak{m}^{3}$ and the cokernel of
  $\left(\begin{smallmatrix}x_{1}&x_{2}&0\\0&x_{1}&x_{2}\end{smallmatrix}\right)$,
  and $k\oplus R/\mathfrak{m}^{2}$; in three variables the residue field,
  $R/\mathfrak{m}^{2}$, two monomial complete intersections, a monomial
  ideal of colength five and the Gorenstein algebra
  $R/(y_{1}^{2}-y_{2}^{2},y_{2}^{2}-y_{3}^{2},y_{1}y_{2},y_{2}y_{3},y_{1}y_{3})$;
  and three complexes with several cohomology modules. In every case the
  Euler characteristic is nonzero and every product of two independent jet
  classes is nonzero, which is \Cref{lem:finiteproduct}. It then builds the
  cone $C$ of $J(\partial_{2})J(\partial_{1})\colon K\to K[2]$ on the Koszul
  complex $K$ of the residue field in two variables, checks that $C$ is a
  complex with cohomology of lengths $1$ and $1$ and Euler characteristic
  zero, and that $J(\partial_{2})J(\partial_{1})$ is null homotopic on $C$,
  while on $k\oplus k[1]$, also of Euler characteristic zero, it is not. That
  is the sharpness of the lemma. The transcript
  \texttt{m2/finite\_\allowbreak length\_\allowbreak products.txt} is
  distributed with the script.

\item \emph{The two classes on a Mumford fourfold, and the K3 surface that
  would carry them.} The script \texttt{mumford\_\allowbreak rm.py} works on
  the split model $V=V_{1}\otimes V_{2}\otimes V_{3}$ of
  \Cref{setup:mumford}, with $\mathfrak{sl}_{2}^{3}$ acting through its nine
  generators and $\psi=\epsilon\otimes\epsilon\otimes\epsilon$, in five parts.
  (A) The commutant of $\mathfrak{sl}_{2}^{3}$ in $\End\bigwedge^{2}V$,
  computed as the null space of a linear system in $784$ unknowns, has
  dimension four, and the four projections built from the three Casimir
  operators are orthogonal idempotents of ranks $1$, $9$, $9$, $9$ that sum to
  the identity, commute with the action and span the commutant. (B) The Hodge
  numbers are $(6,16,6)$ on $H^{2}$, $(0,9,0)$ on $U_{12}$, $(3,3,3)$ on
  $U_{13}$ and $U_{23}$, and $(1,7,1)$ on $T(-1)$. (C) The product map $\mu$ of
  \Cref{lem:mumfordmu} carries $T_{i}\otimes T_{j}$ onto $U_{ij}$ with rank nine
  for $i\ne j$ and $T_{i}\otimes T_{i}$ onto the line of $\theta$, and its
  image is all of $H^{2}$. (D) The Lefschetz pairing
  $\int\alpha\beta\theta^{2}$ is nondegenerate on $H^{2}$; $\mu\mu^{\dagger}$
  acts on the three $U_{ij}$ by one and the same nonzero rational number
  $\nu$, which is $-1$ in the normalisation of the script; for three choices of
  the conjugates of $e$ and $\lambda$, $\mu_{e}\mu_{e}^{\dagger}$ acts on
  $U_{ij}$ by $\nu\sigma_{i}(g)\sigma_{j}(g)$ with $g=e^{2}/\lambda$ and
  preserves the line of $\theta$; the algebra it generates on $U$ is spanned by
  the three projections when the conjugates of $g$ are distinct and is the
  scalars when they coincide; and $\kappa^{\dagger}\kappa$, for the trace form
  on $\End(V)$, is multiplication by $4e^{2}/\lambda$. (E) In the Clifford
  algebra of $T$ with the form $\sum_{i}l_{i}\operatorname{tr}(x_{i}y_{i})$,
  $(l_{1},l_{2},l_{3})=(1,2,5)$, built on bit masks of an orthogonal basis,
  the spin lifts of the three copies of $\mathfrak{sl}_{2}$ act on $C^{+}(T)$,
  of dimension $256$; the vectors of highest weight $(1,1,1)$ span $32$
  dimensions and lowering them gives a basis, so $C^{+}(T)\cong V^{\oplus32}$.
  In that basis the map $\Psi(v)x=vxv_{0}$ of \Cref{prop:ksembed} is
  $v\otimes N_{i}$ for $v\in T_{i}$, the three matrices $N_{i}$ anticommute
  pairwise, square to the scalars $4$, $8$, $20$, proportional to the $l_{i}$,
  and are linearly independent. That is \Cref{prop:mumfordrm},
  \Cref{lem:mumfordmu} and the shape of the Kuga--Satake map used in the proof
  of \Cref{thm:mumfordks}.

\item \emph{The Lefschetz operator of an abelian scheme over a curve.} The
  script \texttt{lefschetz\_\allowbreak family.py} works in exact exterior
  algebra over $\QQ$ and assumes no sign convention. (A) On ten abelian
  varieties with $g\le4$ and polarisation types from $(1)$ to $(1,2,4)$ and
  $(1,1,2,2)$, the contraction by the dual bivector satisfies
  $[L,\Lambda]=H$, and the Pontryagin product, computed from
  $x\star y=\mu_{*}(p_{1}^{*}x\cup p_{2}^{*}y)$ with $\mu_{*}$ defined against
  $\mu^{*}$ by the projection formula, gives $D^{-1}(x\star\gamma)=\Lambda x$
  in every degree, with $\ell\star\gamma=gD$; without the factor $D^{-1}$ the
  identity fails at type $(1,2)$. (B) On four product families $E\times A$
  over an elliptic curve, where every class is invariant, the operator
  $D^{-1}T+\Lambda_{C}$ of the proof of \Cref{thm:lefschetzfamily}, with
  $\Lambda_{C}$ built from bases of invariant classes and the inverse of their
  pairing matrix, satisfies $[L_{W},\Lambda]=H$, and the cross relations
  $[L_{\mathrm{rel}},\Lambda_{C}]=0$ and $[L_{C},T]=0$ hold. (C) The invariants
  of $\mathfrak{sl}_{2}^{3}$ in $\bigwedge^{q}(V_{1}\otimes V_{2}\otimes V_{3})$,
  computed as the joint kernel of the six generators, have dimensions
  $1,0,1,0,1,0,1,0,1$. That is \Cref{prop:pontryaginlambda}, the algebra of
  \Cref{thm:lefschetzfamily} and the count in \Cref{cor:mumfordlefschetz}.

\item \emph{The two remaining targets, reduced.} The script
  \texttt{targets\_\allowbreak reduction.py} works in the split model
  $H^{*}(X\times X)=\bigwedge^{*}(V\oplus V)$, sixteen generators, with
  exact rational arithmetic where a space is determined and ranks modulo the
  prime $1000003$ of integer vectors where a span is to be shown full, a rank
  modulo a prime being a lower bound for the rank over $\QQ$. (A) The
  invariants of $\mathfrak{sl}_{2}^{3}$ in $\bigwedge^{k}(V\oplus V)$ have
  dimensions $1,0,3,0,8,0,16,0,28,0,16,0,8,0,3,0,1$; in degree four the
  products of divisor classes give six, so two are exceptional; and the ring
  of invariants is generated in degrees two and four. (B) At a CM point the
  Hodge classes of $X_{c}$ number $1,0,4,0,8,0,4,0,1$ and those of
  $X_{c}\times X_{c}$ number $1,0,16,0,132,0,432,0,648,\dots$; the divisor
  classes of $X_{c}\times X_{c}$ span only $100$ of the $132$ in degree four,
  while the ring generated by them and by the pull-backs of the Hodge classes of
  $X_{c}$ along $(\varphi,\psi)\colon X_{c}\times X_{c}\to X_{c}$ spans every
  Hodge class in every degree. (C) On an abelian $2n$-fold of Weil type,
  $n=1,2,3$, the classes of type $(p,p)$ annihilated by $Q\wedge HH^{1}$, by
  $P\wedge HH^{1}$ and by $HH^{1}\wedge HH^{1}$ are $\CC\alpha_{-}$,
  $\CC\alpha_{+}$ and $0$. That is \Cref{thm:mumfordcm}, the ring structure
  used in \Cref{cor:mumfordequiv}, and the annihilators in the proof of
  \Cref{thm:nosum}.

\item \emph{The exceptional classes of a Mumford square are rigid.} The
  script \texttt{mumford\_\allowbreak rigidity.py} works in the model of
  \Cref{setup:omegaa} at a point of the Mumford curve, with $V^{1,0}$ the
  vectors of third weight $+1$, and in exact rational arithmetic with the
  coefficients $a_{0},\dots,a_{3}$ as symbols. (A) The Casimir operators
  split ${\textstyle\bigwedge^{2}}V$ as $1+9+9+9$ and the four tensors $\pi$ are invariant.
  (B) $\mathfrak{sp}^{-1,1}$ has dimension $36$ and splits into nine weight
  blocks of dimensions $3,4,3,4,8,4,3,4,3$; the generic kernel of contraction
  into $\omega_{a}$ is one vector, independent of $a$; every Gram determinant
  is a positive constant times a product of $P_{1},P_{2},Q_{1},\dots,Q_{4},R$;
  each is a sum of squares whose zeros force two of $a_{1},a_{2},a_{3}$ to
  coincide; the tangent dimensions at $(3,5,-7,11)$, $(3,5,5,5)$ and
  $(1,1,1,1)$ are $1$, $10$ and $20$. (C) $HT^{2}=28+64+28$ splits into $46$
  blocks by type, weight and exchange parity, $\omega_{a}$ being exchange
  invariant; every Gram determinant is a positive constant times a product of
  the forms of (B), $S_{1},\dots,S_{4}$, $Z_{1},\dots,Z_{6}$, $a_{2}$, $a_{3}$
  and four further linear forms; the only block with a kernel at a generic
  point is the even block of weight zero in $H^{1}(T)$, and its kernel is the
  lowering derivation of the third factor on both copies; $L_{0}$ divides only
  the nine odd bivector blocks. (D) Certificates: each $S$ is a sum of squares
  whose zeros force a coincidence or a zero among $a_{1},a_{2},a_{3}$, each
  $Z$ is a positive combination of four independent squares, and the linear
  factors other than $L_{0}$ have unequal coefficients on $a_{1},a_{2},a_{3}$.
  (E) As a cross-check, the annihilator in $HT^{2}$ has dimension one at forty
  random admissible points, $36$ at $(1,1,1,1)$, $7$ at $a_{2}=0$, $4$ on the
  hyperplane $a_{0}-3a_{1}-3a_{2}+a_{3}=0$, and $10$ at three points of
  $L_{0}=0$. That is \Cref{thm:mumfordrigid}, \Cref{rem:rigidcompare} and
  the annihilator used in \Cref{thm:mumfordnumerical}.

\item \emph{What an object with an exceptional Chern character must look
  like.} The script \texttt{mumford\_\allowbreak object.py} works in the same
  model. (A) At $a=(3,5,-7,11)$ the ranks of contraction into $\omega_{a}$ on
  $HT^{k}$, computed exactly over $\QQ$ block by block, are
  $1,16,119,328,560,328,119,16,1$ for $k=0,\dots,8$ and $0$ above. (B) The
  same ranks hold modulo the prime $1000133$, at which the minimal polynomial
  of $\zeta_{7}+\zeta_{7}^{-1}$ splits, at $a_{0}=2$ and the three conjugates
  in each of the six orders; at a point of $L_{0}=0$ they are
  $1,16,110,328,548,328,110,16,1$. (C) For a random class with components of
  types $(0,0)$ to $(3,3)$ the ranks are palindromic and vanish above degree
  eight. (D) The alternating sum of the ranks is $112$ and
  $2\cdot1+2\cdot119+560=800$. (E) At a CM point there are $16$ divisor
  classes and $132$ Hodge classes in degree four; products of divisor classes
  span $100$, and meet the span of the four $\pi$ exactly where
  $a_{1}=a_{2}=a_{3}$. (F) At a CM point $X_{c}$ has $8$ Hodge classes in
  degree four, of which products of divisor classes span $6$. That is
  \Cref{lem:rankduality}, \Cref{thm:mumfordprofile} and
  \Cref{prop:mumforddivisor}.

\item \emph{What line bundles generate on a Mumford square.} The script
  \texttt{lefschetz\_\allowbreak closure.py} works in the same model, with
  $\mathfrak{sp}(V,\psi)$ spanned by the $36$ maps $x\mapsto\psi(u,x)v+
  \psi(v,x)u$ and the torus of diagonal matrices in $\Sp(V,\psi)$, of rank
  four. (A) The $36$ maps and the Lie algebra of $G$ lie in
  $\mathfrak{sp}(V,\psi)$. (B) The $\Sp(V,\psi)$-invariants of
  ${\textstyle\bigwedge^{k}}(V\oplus V)$ have dimension at most $1,3,6,10,15,10,6,3,1$ for
  $k=0,2,\dots,16$, as a kernel modulo $1000003$, and products of the three
  invariant divisor classes span spaces of exactly these dimensions. (C) The
  $G$-invariants have dimension $1,3,8,16,28,16,8,3,1$, by weight
  multiplicities and as a kernel, so the Hodge classes that are not
  Lefschetz classes number $0,0,2,6,13,6,2,0,0$, twenty-nine in all. (D) At
  a CM point the monomials of weight zero for the diagonal torus number
  $1,16,100,304,454,304,100,16,1$, the coefficients of $(1+4y+y^{2})^{4}$;
  each is a product of the $16$ divisor classes, and those products span
  exactly them; the Hodge classes number $1,16,132,432,648,432,132,16,1$. (E)
  $\pi_{0}$ and $\pi_{12}+\pi_{13}+\pi_{23}$ are $\Sp(V,\psi)$-invariant and
  $\pi_{12}$ is not; the components of $\pi_{12},\pi_{13},\pi_{23}$ of
  nonzero weight for the diagonal torus have rank two over $\QQ$, with the
  relation $(1,1,1)$ only. (F) The modules generated by $\pi_{12}-\pi_{13}$
  and by $\pi_{13}-\pi_{23}$ have dimension $308$ and highest weight
  $(2,2,0,0)$, and Weyl's formula gives $1,27,42,308$ for the highest weights
  $0,\varpi_{2},\varpi_{4},2\varpi_{2}$, with sum $378$. (G) The module meets
  the span of the four $\pi$ in a plane. That is
  \Cref{thm:lefschetzclosure}.

\item \emph{The Hodge locus of an exceptional class among all complex tori.}
  The script \texttt{twistor\_\allowbreak locus.py} realises $V_{\RR}$ as the
  fixed space in $V_{1}\otimes V_{2}\otimes V_{3}$ of the real structure that
  is quaternionic on the first two factors and real on the third, in exact
  Gaussian integer arithmetic. (A) It is a real structure, $V_{\RR}$ has
  dimension $8$, and $\psi$ is real and alternating on it. (B) The complex
  structures $1\otimes1\otimes j_{3}$ of the Mumford family and
  $j\otimes1\otimes1$, $1\otimes j\otimes1$ of the compact factors preserve
  $V_{\RR}$, and $\psi(x,Jy)$ has signature $(0,8)$ for the first and $(4,4)$
  for the other two. The unit quaternion $k\otimes1\otimes1$ with $kj=-jk$
  preserves $V_{\RR}$ and $\psi$ and carries $\psi(x,Jy)$, for
  $J=j\otimes1\otimes1$, to its negative, so every form $\psi\otimes s$ on
  $V_{\RR}\oplus V_{\RR}$ with $s\ne0$ is congruent to its negative and
  indefinite. (C) At $j\otimes1\otimes1$
  the tensors $\pi$ have type $(2,2)$ and the annihilator of $\omega_{a}$ in
  $H^{1}(T)$, of dimension $64$, is spanned by the lowering derivation of the
  first factor; for the Mumford family it is spanned by that of the third,
  as in item (XLIV). (D) The exchange of the first and third factors carries
  $\pi_{12}$ to $\pi_{23}$ and fixes $\pi_{0}$ and $\pi_{13}$. That is
  \Cref{prop:notwistor}.

\item \emph{The square as a holomorphic symplectic variety.} The script
  \texttt{hk\_\allowbreak pullback.py} works in the same split model, with
  $\iota(x)=(x\otimes1)\Psi$ for $x$ in $T=\bigoplus_{i}\mathfrak{sl}(V_{i})$.
  (A) $\iota(x)$ is symmetric and $\iota$ is equivariant, for all pairs of the
  nine generators. (B) The three Casimir operators split
  $\bigwedge^{2}(V\oplus V)$, of dimension $120$, into isotypic pieces of
  highest weights $(0,0,0)$, $(2,2,0)$, $(2,0,2)$, $(0,2,2)$, $(2,2,2)$,
  $(2,0,0)$, $(0,2,0)$, $(0,0,2)$ and dimensions $3,27,27,27,27,3,3,3$, with
  parts of type $(2,0)$ of dimensions $0,0,9,9,9,0,0,1$; the last three pieces
  are spanned by $\iota(T_{1})$, $\iota(T_{2})$, $\iota(T_{3})$. (C) The form
  $\sigma_{0}=\iota(E_{3})$ is of type $(2,0)$ and pairs the vectors of type
  $(1,0)$ of the two copies through an invertible $4\times4$ block. (D)
  $\iota_{2}(C_{1})=3\pi_{0}-\pi_{12}-\pi_{13}+3\pi_{23}$ and cyclically,
  exactly, and hence the formula for $\iota_{2}(\sum s_{i}C_{i})$ with symbolic
  $s_{i}$. (E) $\det(x_{1}\otimes1\otimes1+1\otimes x_{2}\otimes1+
  1\otimes1\otimes x_{3})=\Delta(N_{1},N_{2},N_{3})^{2}$ as a polynomial
  identity in the nine coordinates. (F) $\iota(x)^{8}=8!\det(x)$ times the
  volume form at three random integral points. (G) $\Delta$ is a
  nondegenerate ternary form, of Gram determinant $-4$, and irreducible. That
  is \Cref{prop:hksquare}.

\item \emph{Which of the open routes to the Mumford target are needed.} The
  script \texttt{mumford\_\allowbreak routes.py} carries the implications into
  and out of the Mumford target as data, in the manner of item (XXXIII):
  sixteen statements, one of them proved (the Hodge conjecture at the CM
  points), and twenty-two Horn rules, each labelled by the theorem that proves
  it; unlike that rule set this one contains equivalences. (A) Its eleven
  labels are labels of the paper and its statements are all used. (B) The target is not
  in the closure of what is proved. (C) The statements equivalent to the
  target over what is proved are exactly the target, algebraicity at
  uncountably many points, $B(W\times_{C}W)$, bounded data at infinitely many
  points, and bounded data modulo $p$. (D) The minimal sets of open statements
  that yield the target, found by exhaustive search, number twelve and each
  has one element. (E) The target yields neither the complex with
  $\dim\Ext^{2}=119$ nor the two Kuga--Satake statements. (F) \textup{(L)}
  and \textup{(V)} each yield the target. (G) The Hodge conjecture and (F2)
  yield the target, and the target yields none of the Hodge conjecture, (F1),
  (F2), (F3), \textup{(L)}, \textup{(V)}, \textup{(M)}. (H) (F3) is
  equivalent to the Hodge conjecture by \Cref{prop:f3ishc}, so it is the
  twelfth one-element route to the target, and the minimal sets of open
  statements that yield the Hodge conjecture in this rule set are the
  conjecture itself, (F3), and \textup{(L)} with \textup{(M)}, none of which
  contains the target. That is \Cref{cor:mumfordone} and
  \Cref{rem:mumfordneeded}.
\item \emph{What an object meeting the numerical criterion must look like.}
  The script \texttt{criterion\_\allowbreak shape.py} checks the finite
  ingredients of \Cref{thm:p2endo}. (A) In the coordinate model of
  \Cref{sec:explicit}, for $n=2$ with $d=1,2,3,5,7$ and for $n=3$ with
  $d=1,2,3$, the rational Weil classes $\omega_{1}$, $\omega_{2}$ satisfy
  $\eta\wedge\omega_{1}=\eta\wedge\omega_{2}=0$, and in the orientation in
  which $\eta^{2n}$ is positive the intersection form on the rational Weil
  plane is positive definite at $n=2$ and negative definite at $n=3$, as the
  Hodge--Riemann relations require; its Gram matrix in the basis
  $\omega_{1},\omega_{2}$ is $\mathrm{diag}(8d^{2},8d)$ at $n=2$ and
  $\mathrm{diag}(-32d^{3},-32d^{2})$ at $n=3$. (B) In the model of item
  (XXXVIII), for $n=2$ and $n=3$ and in every degree, the classes of type
  $(p,p)$ killed by all $4n^{2}$ monomials of $P\wedge Q$ are exactly
  $\CC\alpha_{+}\oplus\CC\alpha_{-}$. (C) The generators of
  $\bigwedge^{2n}P$ and $\bigwedge^{2n}Q$ contract $\omega$ to nonzero
  multiples of the generator of $H^{0,2n}$; contraction with $\omega$ is
  injective on $\bigwedge^{2n-2}P\oplus\bigwedge^{2n-2}Q$; and the pairing
  $\bigwedge^{2}P\times\bigwedge^{2n-2}P\to\bigwedge^{2n}P$ is perfect.
  (D) The alternating sum of the lower bounds $1,2\binom{2n}{1},\dots,1$ is
  $-2$ for every $n\le60$; at $n=2$ and $n=3$, exhaustive search over the
  admissible profiles shows $2e_{0}\ge\chi+4$, hence $e_{0}\ge3$, lists the
  smallest profiles, and gives $e_{1}\le e_{0}+9$ at $n=3$; at $n=4$ and
  $n=5$ a profile with $e_{0}=1$ satisfies the same constraints, so the bound
  on $\End$ is special to $n\le3$; and the same identity with $\chi=0$ on the
  Mumford square returns the $800$ odd self-extensions of item (XLV).
\item \emph{The Hodge classes of a very general Weil torus.} The script
  \texttt{weil\_\allowbreak tori.py} checks the finite inputs of
  \Cref{lem:weiltorushodge}, \Cref{lem:weiltorussubsets} and
  \Cref{thm:p2false}, in the coordinate model of \Cref{sec:explicit}, for
  $n=2$ with $d=1,3$ and for $n=3$ with $d=2$. The members of the family $S$
  of \Cref{setup:weiltori} used are $T_{X,Y}$ with $X$ and $Y$ spanned by
  explicit $K$-rational combinations of the eigenvectors, three of them for
  $n=2$ and five for $n=3$. (A) The rational Weil classes $\omega_{1}$,
  $\omega_{2}$ are of type $(n,n)$ at every member used. (B) The
  polarisation $\eta$ is of type $(1,1)$ at the balanced member and at none of
  the members used, which therefore lie off the polarised family. (C) The
  rational classes of degree $2k$, $1\le k\le n$, that are of type $(k,k)$
  at all the members used form a space of dimension $0$ for $k<n$ and $2$ for
  $k=n$, the Weil plane. A class of degree $2k$ is of type $(k,k)$ exactly
  when its product with every product of $2n-k+1$ holomorphic $1$-forms
  vanishes; the ranks are computed modulo a prime $q<2^{31}$ in which $-d$ is
  a square, with $\sqrt{-d}$ sent to both square roots so that the conjugate
  conditions are imposed as well, and since reduction can only lower a rank
  the dimensions found are upper bounds, attained by the Weil classes. Degrees above
  $2n$ are not checked; the proof of the lemma covers them. (D) For an
  element $\kappa$ of $V_{+}\otimes V_{-}$ with random coefficients,
  $\alpha_{+}\wedge\kappa^{n}=\alpha_{-}\wedge\kappa^{n}=0$, and both
  products become nonzero once a component $\beta\in\bigwedge^{2}V_{-}$
  with $\beta^{n}\ne0$ is added, so the vanishing used for the block
  diagonal K\"ahler class $\kappa_{0}$ is a property of $V_{+}\otimes V_{-}$
  and not of the coordinates.
\end{enumerate}

Items (I), (II), (III), (VI), (VII), (VIII) and (X) to (LI) are the ones
on which theorems depend. Items (IV) and (V) support statements that no
obstruction exists, and a failure there would strengthen rather than weaken
the paper, so they are included for completeness rather than for safety.
Running \texttt{verify\_all.py} executes the forty-six Python items and
prints \texttt{865 checks passed, 0 failed}; items (XXIX), (XXXI) and
(XL) need Macaulay2 and are run separately, and their transcripts are
distributed with the scripts.
